%% kapa: exam=1 round=19 year=2008 period=1 form=A subject=civil_law subject_ko=민법 seq=25 number=25
%% source: Q-Net 감정평가사 기출문제 (cst00309) artlSeq=1001962 file=extracted/1008/1001962/2008년 제19회 감정평가사 1차시험 기출문제/1교시.hwp
%% conversion: kordoc 4.24.51 -> tools/gen1.py (auto); review=proofread-sonnet; flags=none
\begin{question}{25}
등기에 관한 설명으로 틀린 것은? (다툼이 있는 경우에는 판례에 의함)
\begin{choices}
  \item 위조문서에 의한 등기는 그것이 실체관계에 부합하거나 그 등기에 부합하는 물권행위가 있는 한 유효하다.
  \item 소유권보존등기명의인이 동일인이 아닌 경우에는 비록 나중에 이루어진 보존등기라 할지라도 그것이 실체관계에 부합하면 유효하다.
  \item 중간생략등기의 합의 하에 최종 매수인과 최초 매도인을 당사자로 하는 토지거래허가를 받아 최초 매도인으로부터 최종 매수인 앞으로 경료된 소유권이전등기는 무효이다.
  \item 등기명의인을 달리하여 중복등기가 된 경우에는 원칙적으로 제2등기가 무효이지만, 예외적으로 제1등기가 원인무효이고 제2등기가 실체관계에 부합하는 경우에는 제2등기가 유효하다.
  \item 동일인명의로 소유권보존등기가 중복등기된 경우, 그 실체관계의 부합 여부를 가릴 것 없이 먼저 행하여진 등기가 유효하다.
\end{choices}
\end{question}
\answer{2}
